\bookchapter{บทที่ 26:แนวโน้มเทคโนโลยีโทรศัพท์ในปัจจุบัน}

\chapterintro{พูดคุยเรื่อง “ชีวิตหลังเกษียณ” ของโทรศัพท์เก่ากันไปแล้ว ตอนนี้ลองหันกลับมามองปัจจุบันกันบ้าง ปี 2024-2026 ถือเป็นช่วงที่เทคโนโลยีโทรศัพท์มือถือเติบโตแบบก้าวกระโดด: AI ถูกยัดลงมาอยู่ในกระเป๋ากางเกง, การสื่อสารผ่านดาวเทียมเริ่มจากเครื่องระดับ Flagship แล้วลงมาสู่รุ่นระดับกลาง, หน้าจอแบบพับได้จาก “ของลองใหม่” กลายเป็น “ของใช้ประจำวัน” และเทคโนโลยีแบตเตอรี่ก็มีความก้าวหน้าระดับวัสดุที่ไม่ได้เห็นมานานหลายปี ในบทนี้ เราจะมาสรุปแนวโน้มเทคโนโลยีล้ำสมัยที่กำลังเกิดขึ้นเหล่านี้ เพื่อช่วยให้คุณตัดสินใจว่า “ถึงเวลาเปลี่ยนเครื่องหรือยัง” และ “ควรเปลี่ยนเป็นเครื่องแบบไหน”}

\booksection{โทรศัพท์ AI: เป็นแค่เครื่องมือทางการตลาด หรือการเปลี่ยนแปลงที่แท้จริง?}

ตั้งแต่ปี 2024 เป็นต้นมา ผู้ผลิต Android แทบทุกรายต่างผลักดันแนวคิด “โทรศัพท์ AI” บนเครื่องระดับ Flagship และ Apple ก็เปิดตัว Apple Intelligence ใน iOS 18 แต่จริงๆ แล้วโทรศัพท์แบบไหนถึงจะถูกเรียกว่า “โทรศัพท์ AI”? ในวงการมีมาตรฐานในการพิจารณาอยู่ประมาณ 3 ข้อ:

\begin{itemize}
  \item \textbf{Large Model บนอุปกรณ์}: โทรศัพท์สามารถรันโมเดลภาษาที่มีพารามิเตอร์ประมาณ 7B ได้ในเครื่อง (เช่น เวอร์ชันประหยัดพื้นที่ของ Qwen, Llama, DeepSeek) โดยไม่ต้องพึ่งพาคลาวด์ก็สามารถสรุปข้อความ, แปลภาษา, หรือสร้างข้อความได้ วิธีตรวจสอบ: หากใช้งานฟีเจอร์ AI ได้ในโหมดการบิน แสดงว่าเป็นโมเดลแบบ On-device
  \item \textbf{กำลังประมวลผล NPU (TOPS)}: กำลังประมวลผลของ NPU (หน่วยประมวลผลโครงข่ายประสาทเทียม) มีหน่วยเป็น TOPS (ล้านล้านการดำเนินการต่อวินาที) ชิประดับ Flagship ปี 2024 (เช่น Snapdragon 8 Gen3, Dimensity 9300) มีกำลังประมวลผล NPU อยู่ในช่วง 40-70 TOPS ส่วนชิปรุ่นใหม่ในปี 2025-2026 จะทะลุ 100 TOPS กันเป็นวงเลิก ซึ่งนี่คือสเปคฮาร์ดแวร์ขั้นต่ำในการรันโมเดลภาษาขนาดใหญ่บนเครื่อง
  \item \textbf{สถานการณ์การใช้งาน AI Agent}: ไม่ใช่แค่ตอบคำถามได้เท่านั้น แต่ยังสามารถทำงานเป็นขั้นตอนข้ามแอปได้ด้วย — เช่น "ช่วยจัดรูปถ่ายจากการประชุมเมื่อวานให้เป็นอัลบั้มแล้วแชร์ให้ Zhang San" ระบบจะดึงแอปต่างๆ เช่น อัลบั้ม, ไฟล์, WeChat มาทำงานร่วมกันอัตโนมัติ นี่คือสัญลักษณ์สำคัญที่ทำให้โทรศัพท์ AI แตกต่างจาก "โทรศัพท์ที่มีผู้ช่วย AI"
\end{itemize}

\booksection{เปรียบเทียบกลยุทธ์ AI ของแต่ละแบรนด์: ใครทำจริง ใครแค่เอาชื่อไปแปะ?}

เส้นทาง AI ของแต่ละแบรนด์มีความแตกต่างกันอย่างชัดเจน เวลาเลือกซื้อสามารถโฟกัสที่จุดเหล่านี้ได้:

\begin{itemize}
  \item \textbf{Apple Intelligence (iOS 18+)}: เน้นการผสานรวมกับแอประดับระบบ — สรุปอีเมล, จัดลำดับความสำคัญของการแจ้งเตือน, ช่วยเขียนและปรับแก้ข้อความ, และแก้ไขรูปภาพด้วย AI ในอัลบั้ม โดยใช้โมเดล On-device เป็นหลัก แต่หากเป็นงานที่ซับซ้อนจะส่งขึ้นไปประมวลผลบน Private Cloud Compute จุดแข็งคือการปกป้องความเป็นส่วนตัวและความลื่นไหลของระบบ ส่วนจุดอ่อนคือรองรับภาษาจีนช้าไปหน่อย บางฟีเจอร์จะค่อยๆ รองรับในปี 2025
  \item \textbf{Google Gemini Nano}: ฝังมาในซีรีส์ Pixel 8/9 รองรับ "วงกลมเพื่อค้นหา", Magic Compose, สรุปไฟล์เสียงบันทึก ฯลฯ Gemini Nano เป็นโมเดลภาษาขนาดเล็กแบบ On-device ที่แท้จริง เป็นมิตรกับระบบนิเวศโอเพนซอร์ส และถือเป็นมาตรฐานอ้างอิงสำหรับโทรศัพท์ AI ของ Android
  \item \textbf{โมเดลใหญ่ Pangu ของ Huawei}: เชื่อมต่อกับ XiaoYi ของ HarmonyOS รองรับการสร้างข้อความ On-device, แก้ไขรูปภาพด้วย AI, ลบวัตถุที่ไม่ต้องการด้วย AI (เช่น คนเดินผ่าน), และสรุปการโทร หลังจาก HarmonyOS NEXT 5.0 จะมีการผสานรวมระดับระบบที่ลึกซึ้งยิ่งขึ้น โดยความสามารถด้าน AI จะหลอมรวมกับระบบปฏิบัติการอย่างแนบแน่น
  \item \textbf{โมเดลภาษาขนาดใหญ่ MiLM ของ Xiaomi}: ฝังมาใน HyperOS รองรับ AI Portrait, สรุปข้อความด้วย AI, AI Painter, และ Super Resolution โดย Xiaomi ใช้เส้นทาง "ประสานงานระหว่างเครื่องและคลาวด์" คือรันโมเดลขนาดเล็กบนเครื่อง และส่งงานหนักขึ้นคลาวด์
  \item \textbf{OPPO (โมเดล Andes) และ vivo (โมเดล BlueLM)}: ผสานรวมอย่างลึกซึ้งใน ColorOS 15 และ OriginOS 5 โดยเน้นฟีเจอร์ AI ลบวัตถุ, สรุปเอกสารด้วย AI, ถามตอบเกี่ยวกับรูปภาพด้วย AI, และค้นหารูปภาพด้วยภาษาธรรมชาติ ("หารูปที่ถ่ายริมทะเลเมื่อปีที่แล้ว")
\end{itemize}

\textbf{คำแนะนำในการเลือกซื้อ}: หากคุณต้องใช้ AI จัดการเอกสารงานและประมวลผลข้อมูลบ่อยๆ ให้เลือกเครื่องที่มีกำลังประมวลผล On-device ที่แข็งแกร่งและมีการผสานรวม AI ระดับระบบ แต่ถ้าแค่ใช้แก้ไขรูปภาพด้วย AI บางครั้ง โทรศัพท์ระดับกลางรุ่นหลักๆ ก็เพียงพอแล้ว

\booksection{การสื่อสารผ่านดาวเทียม: จาก “จุดขาย” ของ Flagship สู่มาตรฐานของรุ่นระดับกลาง}

การสื่อสารผ่านดาวเทียมถือเป็นหนึ่งในเทคโนโลยีที่เข้าถึงผู้ใช้งานได้จริงที่สุดในช่วงสองปีที่ผ่านมา ตอนปี 2022 ที่ Huawei Mate 60 และ iPhone 14 เปิดตัว ยังเป็นสิทธิพิเศษเฉพาะเครื่องระดับ Flagship แต่ในปี 2025 ก็ได้ลงมาถึงระดับราคา 2,000-3,000 หยวนแล้ว

\begin{itemize}
  \item \textbf{Beidou Short Message ลงมาสู่รุ่นระดับกลาง}: ฟังก์ชัน Beidou Short Message ของดาวเทียม Beidou (ส่ง SMS ได้ไม่เกิน 1,000 ตัวอักษรจีน + พิกัดตำแหน่ง) ได้ลงมาถึงรุ่นระดับกลางของ Redmi, realme, iQOO และแบรนด์อื่นๆ โดยไม่มีค่ารายเดือน ใช้งานได้ฟรี เหมาะสำหรับการเดินป่าหรือขับรถท่องเที่ยว
  \item \textbf{Tiantong Satellite Call}: ซีรีส์ Huawei Mate 60/70 และ Honor Magic6 รองรับฟังก์ชัน Tiantong-1 Satellite Call สามารถโทรออก/รับสายดาวเทียมและส่งข้อความผ่านดาวเทียมได้ โดยต้องสมัครแพ็กเกจดาวเทียมกับผู้ให้บริการ (ประมาณ 10-200 หยวนต่อเดือน) เหมาะสำหรับผู้ที่เข้าไปในพื้นที่ทุรกันดารบ่อยครั้ง
  \item \textbf{iPhone Satellite Messages (SOS)}: ตั้งแต่ iPhone 14 เป็นต้นไปรองรับ Satellite SOS สามารถส่งข้อความไปยังหน่วยงานฉุกเฉินและแชร์พิกัดตำแหน่งได้ โดยให้บริการฟรี 2 ปี และในบางพื้นที่สามารถต่ออายุใช้งานต่อได้
  \item \textbf{คำแนะนำที่ใช้ได้จริง}: ผู้ใช้งานในเมืองแทบไม่ต้องใช้เลย แต่สำหรับนักผจญภัยกลางแจ้งหรือนักธุรกิจที่เดินทางไปพื้นที่ห่างไกล แนะนำอย่างยิ่งให้เลือกเครื่องที่มีฟังก์ชันดาวเทียม — เพราะในยามวิกฤตอาจช่วยชีวิตคุณไว้ได้
\end{itemize}

\booksection{ความพร้อมของระบบนิเวศหน้าจอพับ: ใช้เป็นเครื่องหลักได้หรือยัง?}

หน้าจอแบบพับได้ก้าวจาก “เทคโนโลยีลองใหม่” ไปสู่ “ใช้เป็นเครื่องหลักได้จริง” การจะตัดสินว่าหน้าจอพับพร้อมหรือยัง ให้ดู 3 จุด: ความทนทานของฮาร์ดแวร์, ความเข้ากันของแอป, และระบบนิเวศการโต้ตอบ

\begin{itemize}
  \item \textbf{ความทนทานของฮาร์ดแวร์}: อายุการใช้งานบานพับส่วนใหญ่อยู่ที่ 500,000 ครั้งขึ้นไป (เปิด-ปิดวันละ 100 ครั้ง ใช้งานได้นาน 13 ปี) บานพับแบบหยดน้ำช่วยลดรอยพับบนหน้าจอ และกันน้ำมาตรฐาน IPX8 กลายเป็นมาตรฐานไปแล้ว การใช้กระจก UTG + ฟิล์มป้องกันบนหน้าจอภายใน ทำให้หน้าจอไม่ "บอบบาง" เหมือนเมื่อก่อน
  \item \textbf{ความเข้ากันของแอป}: แอปหลัก (WeChat, Taobao, Douyin, Bilibili, Office, WPS) ได้ปรับแต่งให้รองรับหน้าจอขนาดใหญ่แล้ว แต่ยังมีแอปรองราว 30\% ที่ยังไม่รองรับ ทำให้ภาพถูกยืดหรือมีขอบดำ แนะนำให้ทดลองใช้แอปที่คุณใช้บ่อยก่อนตัดสินใจซื้อ
  \item \textbf{การโต้ตอบในโหมดพับครึ่ง}: เมื่อพับที่มุม 90°-120° แอปบางตัวจะรองรับโหมดพับครึ่งแบบ "ครึ่งบนแสดงเนื้อหา ครึ่งล่างใช้ควบคุม" เหมาะสำหรับการโทรแบบเห็นหน้า, ถ่ายเซลฟี่ และดูวิดีโอ การถ่ายเซลฟีแบบพับครึ่งถือเป็นประสบการณ์ที่มีเฉพาะในโทรศัพท์หน้าจอพับ
  \item \textbf{แบ่งจอและมัลติทาสก์}: บนหน้าจอใหญ่สามารถเปิดแอปแบ่งจอ 2 แอปพร้อมกันได้ หรือแม้กระทั่งเปิดหน้าต่างลอยขนาดเล็ก 3 แอปพร้อมกัน ตรรกะการแบ่งจอของ HyperOS, ColorOS, OriginOS ถือว่าค่อนข้างสมบูรณ์แล้ว แต่สำหรับ iOS (สไตล์ iPadOS) ยังไม่เปิดให้ใช้งานการแบ่งจอบน iPhone
  \item \textbf{ราคา}:แบบพับหน้าจอซ้อน (สไตล์หนังสือ) ลดลงจากระดับหมื่นหยวนมาอยู่ที่ช่วง 5,000-7,000 หยวน ส่วนแบบพับแนวตั้ง (สไตล์ฝาพับ) ลดลงมาที่ช่วง 4,000-5,000 หยวนแล้ว ในปี 2025 ผู้ผลิตบางรายจะเปิดตัวแบบพับแนวตั้งในระดับราคา 3,000 หยวน ทำให้หน้าจอพับเข้าสู่ตลาด mass market อย่างแท้จริง
\end{itemize}

\textbf{คำแนะนำในการเลือกซื้อ}:หากคุณเป็นผู้ใช้งานหนักด้านการอ่าน ประมวลผลเอกสาร หรือวิดีโอ หน้าจอพับแบบซ้อนน่าลองพิจารณา แต่ถ้าให้ความสำคัญกับความพกพาสะดวกและความเป็นสังคม แบบพับแนวตั้งจะเหมาะกว่า ขอแนะนำให้ไปลองจับจับที่ร้านออฟไลน์อย่างน้อย 30 นาทีก่อนตัดสินใจ

\booksection{6G และ Wi-Fi 7: ความเร็วและขีดจำกัดของการเชื่อมต่อยุคถัดไป}

แม้ว่า 6G จะยังไม่มาถึงจริงๆ แต่เทคโนโลยีที่เกี่ยวข้องก็กำลังอยู่ระหว่างการพัฒนาแล้ว:

\begin{itemize}
  \item \textbf{ความคืบหน้าของ 6G}:3GPP วางแผนจะกำหนดมาตรฐาน 6G ให้เสร็จสิ้นภายในปี 2027 และคาดว่าจะเริ่มใช้งานเชิงพาณิชย์ประมาณปี 2030 ความเร็วสูงสุดในทางทฤษฎีของ 6G สามารถสูงถึง 1 Tbps (เร็วกว่า 5G 50 เท่า) มีความหน่วงต่ำเพียง 0.1ms รองรับ "การรวมการสื่อสารและการรับรู้" (Communication + Sensing) จนโทรศัพท์สามารถทำหน้าที่เสมือนเรดาร์ความแม่นยำต่ำได้ ปัจจุบันผู้ผลิตต่างๆ ยังอยู่ในขั้นตอนตรวจสอบต้นแบบ ซึ่งยังห่างไกลจากผลิตภัณฑ์สำหรับผู้บริโภค
  \item \textbf{Wi-Fi 7(802.11be)}:เริ่มเป็นมาตรฐานสำหรับโทรศัพท์เรือธงตั้งแต่ปี 2024 ความเร็วสูงสุดในทางทฤษฎีสูงถึง 46 Gbps (เร็วกว่า Wi-Fi 6 ถึง 4 เท่า) รองรับความกว้างแบนด์ 320MHz, 4K-QAM และ MLO (ส่งสัญญาณหลายลิงก์พร้อมกัน) จากการทดสอบจริงภายใต้เราเตอร์ Wi-Fi 7 โทรศัพท์เรือธงสามารถทำความเร็วส่งถ่ายข้อมูลได้ 2-5 Gbps ซึ่งเร็วกว่า Wi-Fi 6E ถึง 2-3 เท่า หากคุณเปลี่ยนมาใช้โทรศัพท์ที่รองรับ Wi-Fi 7 อย่าลืมจัดหาเราเตอร์ Wi-Fi 7 มาใช้ด้วยจึงจะสามารถใช้ประสิทธิภาพได้เต็มที่
  \item \textbf{Bluetooth 5.4 และ LE Audio}:โทรศัพท์เครื่องใหม่ระหว่างปี 2024-2026 จะรองรับ Bluetooth 5.4 และ LE Audio (เสียงประหยัดพลังงาน) อย่างแพร่หลาย สามารถเชื่อมต่อหูฟัง 2 คู่พร้อมกันเพื่อฟังเพลงเดียวกันได้ ในขณะที่ฟังก์ชัน Auracast สามารถกระจายเสียงให้หลายคนฟังจากแหล่งเสียงเดียวกันพร้อมกันในที่สาธารณะได้
\end{itemize}

\booksection{แบตเตอรี่ซิลิคอน-คาร์บอน: การใช้งานแบตเตอรี่ได้รับการพัฒนาที่ระดับวัสดุในที่สุด}

อายุการใช้งานแบตเตอรี่ของโทรศัพท์มักถูกวิจารณ์มาหลายปีว่าต้องชาร์จทุกวัน แต่ตั้งแต่ปี 2024 เป็นต้นมา การแพร่หลายของเทคโนโลยีแบตเตอรี่ซิลิคอน-คาร์บอน ทำให้สถานการณ์ดีขึ้นอย่างเป็นรูปธรรม

\begin{itemize}
  \item \textbf{หลักการทางเทคโนโลยี}:ความจุทางทฤษฎีของแอโนดแกรไฟต์แบบดั้งเดิมอยู่ที่ประมาณ 372 mAh/g ในขณะที่ซิลิคอนมีความจุทางทฤษฎีสูงถึง 4200 mAh/g (ประมาณ 10 เท่า) แต่ซิลิคอนจะขยายตัวถึง 300\% ขณะชาร์จและคายประจุ ทำให้เกิดการแตกหักได้ง่าย แอโนดซิลิคอน-คาร์บอนใช้โครงสร้างนาโนคอมโพสิตของซิลิคอน-คาร์บอนเพื่อควบคุมปัญหาการขยายตัวของซิลิคอนให้อยู่ในระดับที่ยอมรับได้ ช่วยเพิ่มความหนาแน่นของพลังงานแบตเตอรี่ขึ้น 15\%-30\%
  \item \textbf{ผลการใช้งานจริง}:โทรศัพท์เรือธงที่ใช้แอโนดซิลิคอน-คาร์บอนในช่วงปี 2024-2025 (เช่น Xiaomi 14 Ultra, Honor Magic 6 Ultimate, vivo X100 Pro ฯลฯ) สามารถบรรจุแบตเตอรี่ขนาด 5,000-6,000 mAh ไว้ในตัวเครื่องที่บางเพียง 7-8 มม. ซึ่งมากกว่ารุ่นก่อนหน้าประมาณ 800-1,000 mAh และสามารถใช้งานหนักได้นานถึงวันครึ่ง
  \item \textbf{ทำงานร่วมกับการชาร์จเร็ว 100W}:แอโนดซิลิคอน-คาร์บอนเข้ากันได้ดีกับการชาร์จเร็ว การชาร์จเต็มด้วยความเร็ว 100W/120W ภายใน 17 นาทีไม่ใช่เรื่องใหม่อีกต่อไป ในปี 2025 มีบางรุ่นที่รองรับการชาร์จ 150W-240W สามารถชาร์จถึง 80\% ได้ภายใน 10 นาที
  \item \textbf{คำแนะนำในการเลือกซื้อ}:หากคุณเป็นผู้ใช้งานหนัก ควรเลือกโทรศัพท์ที่ผลิตตั้งแต่ปี 2024 เป็นต้นไปซึ่งใช้ชุดคู่ "ซิลิคอน-คาร์บอน + ชาร์จเร็ว 100W" อาการกังวลเรื่องแบตเตอรี่หมดจะลดลงอย่างมาก
\end{itemize}

\booksection{สรุปบทนี้: จำเทคโนโลยีล้ำสมัยเหล่านี้ให้ได้ในหนึ่งประโยค}

\begin{itemize}
  \item โทรศัพท์ AI 3 มาตรฐาน: โมเดลภาษาขนาดใหญ่บนเครื่อง, พลังการประมวลผล NPU TOPS, AI Agent ข้ามแอป
  \item การสื่อสารผ่านดาวเทียม: ข้อความสั้น BeiDou เริ่มลงสู่รุ่นกลาง ส่วนการโทร Tiantong ยังเป็นเอกสิทธิ์ของเรือธง
  \item หน้าจอพับ: ฮาร์ดแวร์พร้อมแล้ว, แอปพลิเคชันรองรับประมาณ 70\%, การใช้งานแบบพับหยุด/แบ่งหน้าจอคือประสบการณ์ที่มีเอกลักษณ์
  \item 6G ยังอยู่ระหว่างทาง, Wi-Fi 7 มาถึงแล้ว, Bluetooth LE Audio เริ่มแพร่หลาย
  \item แบตเตอรี่ซิลิคอน-คาร์บอน + ชาร์จเร็ว 100W: ความก้าวหน้าครั้งใหญ่ทั้งด้านอายุการใช้งานและการชาร์จ
\end{itemize}

\begin{tcbpractice}{ลองทำดู: สร้าง "รายการความต้องการทางเทคโนโลยี" สำหรับโทรศัพท์เครื่องถัดไปของคุณ}
หยิบกระดาษมาแผ่นหนึ่ง แบ่งเป็น 3 ระดับ “ต้องมี/เป็นจุดเด่น/ไม่ต้องการ” แล้วระบุความต้องการของคุณเกี่ยวกับ AI, ดาวเทียม, หน้าจอพับ, Wi-Fi 7, การชาร์จเร็ว และความจุแบตเตอรี่ วิธีนี้จะทำให้คุณสามารถตัดสินใจได้อย่างรวดเร็วว่า “จุดขาย” ที่พนักงานขายแนะนำในร้านนั้น เป็นสิ่งที่คุณต้องการจริง หรือเป็นเพียงการชวนซื้อ ตัวอย่างเช่น:

\begin{itemize}
  \item ต้องมี: แบตเตอรี่มากกว่า 5,000mAh, ชาร์จเร็วมากกว่า 80W, NFC, กันน้ำ IP68
  \item เป็นจุดเด่น: สรุปความด้วย AI บนเครื่อง, Wi-Fi 7, ข้อความสั้น BeiDou
  \item ไม่ต้องการ: หน้าจอพับ (ฉันคุ้นเคยกับแบบตรง), การโทรผ่านดาวเทียม (ไม่ไปไกล), ชาร์จเร็วขั้นสุด 240W
\end{itemize}

\end{tcbpractice}

\begin{tcbsummary}{สรุปบทนี้}
\begin{itemize}
  \item โทรศัพท์ AI 3 มาตรฐานสำหรับตัดสิน: โมเดลภาษาขนาดใหญ่บนเครื่อง, NPU TOPS, สถานการณ์การใช้งาน AI Agent
  \item กลยุทธ์ AI ของผู้ผลิตต่างๆ: Apple เน้นความเป็นส่วนตัวและการเชื่อมต่อ, Google เปิดซอร์สบนเครื่อง, Huawei/Xiaomi/OPPO/vivo ต่างก็มีจุดเด่นของตัวเอง
  \item การสื่อสารผ่านดาวเทียม: ข้อความสั้น BeiDou ลงสู่รุ่นกลาง ส่วนการโทร Tiantong ยังเป็นเอกสิทธิ์ของเรือธง
  \item หน้าจอพับ: ฮาร์ดแวร์พร้อมแล้ว, แอปพลิเคชันรองรับประมาณ 70\%, การพับหยุดและแบ่งหน้าจอคือประสบการณ์ที่ไม่เหมือนใคร
  \item มาตรฐาน 6G จะกำหนดในปี 2027, Wi-Fi 7 แพร่หลายแล้ว, Bluetooth LE Audio เริ่มเป็นจริง
  \item แบตเตอรี่ซิลิคอน-คาร์บอนทำให้ 5,000-6,000mAh กลายเป็นมาตรฐานของเรือธง อาการกังวลเรื่องแบตเตอรี่ลดลงอย่างมาก
\end{itemize}

\end{tcbsummary}
